\section{Gauge Structure of the Standard Model}
\paragraph*{Step 1: H $\otimes$C gives U(1) $\times$ SU(2)}
\paragraph*{Proposition 7.1 (H $\otimes$C $\to$U(1) $\times$ SU(2)). The biquaternion algebra H $\otimes$R C}
generates exactly the gauge group U(1) $\times$ SU(2).
\paragraph*{Proof. SU(2) from unit quaternions. The unit quaternions {q $\in$H : |q| = 1} $\sim$=}
S3 form a Lie group isomorphic to SU(2). The three generators are Tk := $\sigma$k/2 (half Pauli matrices), satisfying the su(2) relations [Tj, Tk] = i$\epsilon$jklTl (verified numerically to machine precision). U(1) from the complex scalar. The imaginary time unit it $\in$iR generates the abelian group ei$\theta$ $\in$U(1) for $\theta$ $\in$R. This commutes with all SU(2) generators since scalar multiplication commutes with matrix multiplication. Hence U(1) $\times$ SU(2), direct product.
\paragraph*{Step 2: Cl(3) ,$\to$O gives SU(3)}
\paragraph*{Proposition 7.2 (Cl(3) has octonionic dimension). The Clifford algebra Cl(3)}
has basis {1, e1, e2, e3, e12, e13, e23, e123}, eight elements, matching dimR O = 8. The map ei 7$\to$oi (octonion units) defines an embedding Cl(3) ,$\to$O. The three bivectors e12, e13, e23 generate su(2) $\subset$g2 = Aut(O).
\paragraph*{Theorem 7.3 (O $\to$SU(3)). The octonions decompose under SU(3)-action}
as O = C $\oplus$C3,
a singlet plus the fundamental (triplet) representation of SU(3). The triplet C3 carries the three kernels (K1, K2, K3) as the three "colors".
\paragraph*{Proof. The automorphism group of O is the exceptional Lie group G2. Under}
the subgroup SU(3) $\subset$G2, the decomposition O = C $\oplus$C3 holds, where C3 is the fundamental representation of SU(3) (three complex dimensions, i.e. six real dimensions). The eight Gell-Mann generators $\lambda$1, . . . , $\lambda$8 of su(3) act on C3: $\lambda$1,2,3 within the K1-K2 subspace (SU(2) subalgebra), $\lambda$4,5 coupling K1 to K3, $\lambda$6,7 coupling K2 to K3, $\lambda$8 as hypercharge. The identification (K1, K2, K3) = three colors of C3 is direct.
\paragraph*{Step 3: The Standard Model gauge group}
\paragraph*{Theorem 7.4 (Standard Model from the Sphaera biquaternion). The biquaternion ground state qn $\in$H $\otimes$C, extended to the octonionic structure}
C $\otimes$H $\otimes$O, generates the Standard Model gauge group U(1) $\times$ SU(2) $\times$ SU(3) with generators arising from: Source Group Physical force Complex scalar it U(1) Electromagnetism Unit quaternions H SU(2) Weak force Octonion triplet C3 $\subset$O SU(3) Strong force
\paragraph*{Proof. Propositions 7.1--7.2 and Theorem 7.3 establish each factor. The three}
factors are independent (direct product structure) because U(1) commutes with SU(2) (Proposition 7.1) and both commute with the SU(3) action on C3 (the SU(3) generators $\lambda$k act only on the color indices, not on the quaternion structure).
\paragraph*{Remark (The role of K3). The third kernel K3 = e1e2 plays three roles simultaneously, in three different frames:}
- In H $\otimes$C: the photon sector (projection of the Heisenberg compensator,
orthogonal to K1, K2);
- In Cl(3): a bivector (e12, generator of su(2));
- In O: the third color of the SU(3) triplet.
By Theorem 2.1, all three interpretations are related by frame changes that leave $\tau$n and C invariant. K3 is the same object seen from three different frames. This is the quantum-field-theoretic realization of the frameindependence observation: it is invariant under the change of frame.

